% Generated draft: requires mathematical and source review.
\documentclass[11pt]{article}
\usepackage[margin=25mm]{geometry}
\usepackage{fontspec}
\setmainfont{Latin Modern Roman}
\usepackage{amsmath,amssymb,mathtools,unicode-math}
\setmathfont{Latin Modern Math}
\usepackage{enumitem,needspace,xurl,hyperref}
\setlength{\parindent}{0pt}
\setlength{\parskip}{5pt}
\setlength{\emergencystretch}{3em}
\newcommand{\N}{\mathbb{N}}
\newcommand{\Z}{\mathbb{Z}}
\newcommand{\Q}{\mathbb{Q}}
\newcommand{\R}{\mathbb{R}}
\newcommand{\C}{\mathbb{C}}
\newcommand{\F}{\mathbb{F}}
\newcommand{\A}{\mathbb{A}}
\newcommand{\PP}{\mathbb{P}}
\newcommand{\E}{\mathbb{E}}
\newcommand{\eps}{\varepsilon}
\newcommand{\dd}{\,\mathrm{d}}
\begin{document}
\section*{UnsolvedMath ID: 11000125}
\textbf{Source problem code:} catalogue entry
\subsection{Definitions and mathematical statement}
Let \(\Gamma_\infty\) and \(T_\infty\) be the source's stable mapping class and stable Torelli groups. The surface/boundary conventions and stabilization maps are not given in this extract. For an oriented surface bundle \(\pi:E\to B\), the usual positive-index Mumford--Morita--Miller class is \(\kappa_j=\pi_!(e(T_\pi E)^{j+1})\in H^{2j}(B;\Q)\). With compatible stabilization and restriction, the source's Question 3.2 asks
\[
 \kappa_{2i}=0\quad\text{in }H^{4i}(BT_\infty;\Q)\ ?
\]
The intended indices appear to be positive \(i\); the extract does not state the range, and the degree-zero class cannot be silently included in a positive-index stable formulation.

The supplied summary attributes to Morita rational vanishing of the odd-index classes \(\kappa_{2i+1}\) on Torelli and identifies the even-class request with his conjectural vanishing. It mentions later surveys and Sakasai without complete references. Those are supplied literature assessments, not independently verified results here.

The source also suggests, tentatively and with expressed doubt, that the map
\[
 \alpha:BT_\infty\longrightarrow
 F=\operatorname{hofiber}\bigl(B\Gamma_\infty^+
                 \longrightarrow B\mathrm{Sp}(\Z)^+\bigr)\tag{5}
\]
might be zero on cohomology. Here \(B\) denotes classifying space, \(\operatorname{hofiber}\) homotopy fiber, and \(+\) the source's plus-construction convention. The symplectic homology action motivates the map, but the relevant stabilization, plus constructions and lift/nullhomotopy data are missing. Cohomology is contravariant, so a positive-degree version would mean \(\alpha^*:H^{>0}(F;R)\to H^{>0}(BT_\infty;R)\) is zero for the intended coefficients \(R\). Literal zero on all unreduced cohomology cannot hold for nonempty spaces because the unit in degree zero pulls back to the unit. Coefficients and reduced/positive-degree scope therefore require clarification; no full all-degree vanishing theorem is asserted.

Finally, the source reports that [4] finds abundant torsion in \(H^\bullet(B\Gamma_\infty;\Z)\), including classes of every order, but that their description is indirect and complicated. It presents potential wider applications as motivation. This is separate from the rational even-class question and does not imply torsion statements about \(BT_\infty\).
\subsection{Short English statement}
Do the positive even-index MMM classes vanish rationally on the stable Torelli group? The source also raises a tentative stronger cohomological-vanishing question for a map to a homotopy fiber and discusses unrelated stable mapping-class torsion. Its stabilization, coefficients and degree-zero conventions need source review.
\subsection{Sources}
\begin{itemize}
\item Local UnsolvedMath registry, permanent ID \(11000125\), exact source code \texttt{catalogue entry}, Question 3.2, PDF page 104. The supplied title is only ``Question 3.2 --- .'' No individual proposer or proposed year is supplied.
\item Farb, edited volume \emph{Problems on Mapping Class Groups and Related Topics}: \url{https://www.math.uchicago.edu/~farb/papers/mcgbook.pdf}. The volume has 329 explicitly typeset numbered items across chapters, repeated local numbering distinguished by PDF page.
\item Internal reference [4] for stable integral torsion calculations; full entry absent. The research summary names Morita for odd-class vanishing and the even-class conjecture, and Sakasai for related later work, without complete citations.
\item Extraction pdf-font-structure, accessed 2026-07-29. Source evidence NEEDS\_REVIEW presents an open problem and calls for a current primary-literature audit. Supplied OPEN-TRIAGE summary says the even-class and stronger homotopy-fiber questions remain unresolved to its knowledge, with no located stable-range computation settling even-class nonvanishing/vanishing. This is a supplied bounded assessment, not an independent search.
\item Dataset attribution: Ulam AI / UnsolvedMath Contributors, CC BY 4.0. Source rights NEEDS\_REVIEW: publicly accessible author/publisher copy, redistribution terms requiring review.
\end{itemize}
\end{document}
